\subsection{阿廷环上的单模}

设 A 是环。约定俗成地，若左理想 \(\alpha\) 是 A 的非零左理想组成的集合（按包含关系排序）中的一个极小元，我们就称它为 A 的极小左理想。右极小理想与双边极小理想类似地定义。

设 \(\mathfrak{a}\) 是 A 的一个左理想。则 \(\mathfrak{a}\) 是单 A-模当且仅当它是 A 的极小左理想。

左阿廷环（VIII，第 1 页，定义 1）的每个非零左理想都包含一个极小左理想。

命题 4. —— 设 A 是一个具有极小左理想 a 的环。每个忠实单 A-模都同构于 A-模 a。

设 M 是一个忠实单 A-模。设 \(\alpha\) 是 a 的一个非零元素。由于 A-模 M 是忠实的，存在 M 的一个元素 x，使 \(\alpha x \neq 0\) 。于是从 a 到 M 的同态 \(a \mapsto ax\) 非零；从而由 VIII，第 47 页的命题 2，它是同构。

命题 5. —— 设 A 是一个左阿廷环，M 是一个忠实 A-模。

\begin{enumerate}
\def\labelenumi{\alph{enumi})}
\item
  存在一个自然数 \(m\) 与一个从 \(\mathbf{A}_s\) 到 \(\mathbf{M}^m\) 的一个子模的同构。
\item
  若 M 有一个若尔当–赫尔德列 \((\mathrm{M}_{i})_{0\leqslant i\leqslant n}\)，则每个单 A-模都同构于诸商 \(M_{i}/M_{i+1}\) 之一。
\end{enumerate}

把 VIII，第 2 页应用于 A-模 \(A_{s}\) 与 M 的元素的零化子，存在 M 的元素 \(x_{1},\ldots,x_{m}\)，使 M 的零化子是诸 \(x_{i}\) 的零化子之交。由于 A-模 M 是忠实的，A-线性映射 \(a\mapsto(ax_{1},\ldots,ax_{m})\)（从 \(A_{s}\) 到 \(M^{m}\)）是单射，由此得 a)。

在 b) 的假设下，\(A_{s}\) 的每个单商都同构于 \(M^{m}\) 的某个若尔当–赫尔德列的一个商（I, §4, No.~7, 第 43 页，推论），从而同构于诸模 \(M_{i}/M_{i+1}\) 之一（I, §4, No.~7, 第 43 页，定理 6）。我们断定每个单 A-模都同构于 \(A_{s}\) 的一个商。

注. —— 命题 5 特别适用于下面两种情形：

\begin{enumerate}
\def\labelenumi{\alph{enumi})}
\item
  设 A 是交换域 K 上的一个代数，M 是 K 上的一个忠实的有限维模。则 M 是有限长度的 A-模，且 M 的反模是有限生成的。环 A 是左阿廷环（VIII，第 9 页，命题 6）。存在 A-模 M 的一个若尔当–赫尔德列 \((\mathrm{M}_{i})_{0\leqslant i\leqslant n}\)，且每个单 A-模都同构于诸模 \(M_{i}/M_{i+1}\)（对 \(0\leqslant i\leqslant n-1\)）之一。
\item
  设 A 是一个左阿廷环。模 \(A_{s}\) 长度有限（VIII，第 6 页，定理 1）。由于 A-模 \(A_{s}\) 长度有限，存在 A 的左理想的一个递减序列 \((\mathfrak{a}_{i})_{0\leqslant i\leqslant n}\)，使得 \(a_{0}=A\)、\(a_{n}=0\)，且 A-模 \(S_{i}=\mathfrak{a}_{i-1}/\mathfrak{a}_{i}\)（对 \(1\leqslant i\leqslant n\)）都是单模。则每个单 A-模都同构于诸模 \(S_{1},\ldots,S_{n}\) 之一。
\end{enumerate}

